\documentclass{article}

\usepackage[most]{tcolorbox}
\usepackage{physics}
\usepackage{graphicx}
\usepackage{float}
\usepackage{amsmath}
\usepackage{amssymb}


\usepackage[utf8]{inputenc}
\usepackage[a4paper, margin=1in]{geometry} % Controla los márgenes
\usepackage{titling}

\title{Clase Mecanica Estadistica }
\author{Manuel Garcia.}
\date{\today}

\renewcommand{\maketitlehooka}{%
  \centering
  \vspace*{0.05cm} % Espacio vertical antes del título
}

\renewcommand{\maketitlehookd}{%
  \vspace*{2cm} % Espacio vertical después de la fecha
}

\newcommand{\caja}[3]{%
  \begin{tcolorbox}[colback=#1!5!white,colframe=#1!25!black,title=#2]
    #3
  \end{tcolorbox}%
}

\begin{document}
\maketitle

\section{Ejemplo}
Particula libre de masa $ M  $ en el intervalo $ 0<q<1  $ dado que en $ q = 0  $ y $ q = 1  $ hay paredes que ocasionan que la particula rebote (coloicion elastica). 

Luego $ V(q)=0  $; así: 
\begin{gather*}
  H(p,q) = T(p) + V(q) = \frac{p^2 }{2m } = \epsilon 
\end{gather*}

\textbf{a) } solucion de las ecuaciones de movimiento 
\begin{gather*}
  \dot q = \frac{\partial H  }{\partial p } \quad \rightarrow \quad \dot q = \frac{\partial q  }{\partial t } = \frac{\partial  }{\partial p } \left(\frac{p^2 }{2m }\right) = \frac{p }{ m } \quad \rightarrow \quad \displaystyle\int_{q = 0 }^{q } dq = \displaystyle\int_{t = 0 }^{t } \frac{p}{m } dt  \quad \rightarrow \quad q(t) - q(t_0 ) = \frac{p_0 }{m } 
  \\
  \dot p = - \frac{\partial  }{\partial q } ( 0 ) = 0 \quad \rightarrow \quad p(t) = p_0 = m \dot q_0 
  \\
  q(t) = q_0 + \dot q_0 t 
\end{gather*}
\begin{gather*}
  \frac{p_0^2 }{2m } = \epsilon_0 \quad \rightarrow \quad p_0^2 = 2m \epsilon_0  
\end{gather*}

Volumen espacio de fase 
\begin{align*}
  \Lambda(\epsilon_0 ) &= (2 \left|p_0 \right|) (L) \\
          &= 2 \left|p_0 \right|L = 2L \sqrt{2m \epsilon_0 } 
\end{align*}


\section{Ejemplo }
\begin{gather*}
  H(p, q) = \frac{p^2 }{2m } + \frac{1}{2} kq^2    
\end{gather*}
Donde $ V(q) = \frac{1}{2} k q^2  \quad \rightarrow \quad f(q) = - \frac{\partial V  }{\partial q } = - k q   $

\hfill 

\hfill 

\textbf{a) } Slucion ec. de hamilton 
\begin{gather*}
  \dot q = \frac{\partial H  }{\partial p } \quad \rightarrow \quad \dot q = \frac{d q  }{d t } = \frac{p }{m } 
  \\
  \dot p = - \frac{\partial H  }{\partial  q } \quad \rightarrow \quad \dot p = \frac{d p  }{d t } = - k q \quad \rightarrow \quad q = -\frac{1 }{k } \dot p \\
  \dot q = - \frac{1}{k} \ddot p \quad \rightarrow \quad - \frac{1}{k } \ddot p = \dot q = \frac{p }{m }
\end{gather*}
Entonces 
\begin{gather*}
  \ddot p + \frac{k}{m} = 0  \qquad \rightarrow \qquad \ddot p + \omega^2 p = 0 
  \\
  p(t) = B \sin{\omega t } 
\end{gather*}
Sustituyendo 
\begin{gather*}
  \frac{d q  }{d t } = \frac{p}{m } = \frac{B }{m } \sin{\omega t }
  \\
  dq = \frac{B }{m } \sin{\omega t } dt 
  \\
  \displaystyle\int_{q_0 }^{q(t) } dq = \frac{B }{m } \displaystyle\int_{t = 0 }^{t } \sin{\omega t } dt \\
  q(t) = q_0 - \frac{B }{m\omega } \cos{\omega t } + \frac{B }{m \omega} 
\end{gather*}


\caja{blue}{Ejercicio 1 }{
  
}

\end{document}
